\scanpage{8}

10) Soit $M_\alpha$ un système d' \deleted{$\otimes$-générateurs de} \inserted{\unsure{élém$^{\rm ts}$|élém.} \deleted{\unsure{de}} hom. de} $\mathcal M$, tel que avec $K(1)$ et $K(-1)$ ils engendrent $\mathcal M$ comme $\otimes$-catégorie. Alors une polarisation est connue quand, pour tt $\alpha$, on connaît \uline{une} forme de polarisation \inserted{$\varphi_\alpha$} de $M_\alpha$.

On peut supposer d'ailleurs que les $M_\alpha$ sont tous de degré $0$ ou $1$, et si $\mathcal M$ est à engendrement fini, que le syst. est réduit à un $M^0$ et un $M^1$. \deleted{Ils correspondent à des \unsure{=}}

11) \uline{Remarque 1} \deleted{Les polarisati} Pour $M,L$ fixés, les classes d'équivalence de formes de Weil \deleted{\unsure{sur}} $M\otimes M\to L$, \deleted{pour} pour $\mathcal M$ ou $\mathcal M\otimes_K\mathbb R$, sont les $=$. La compatibilité des formes de Weil est la $=$, au sens \deleted{$\mathcal M\otimes_K\mathbb R$} $\mathcal M$ ou $\mathcal M\otimes_K\mathbb R$. Ils correspondent à certaines composantes connexes de l'ouvert de $\operatorname{Hom}(M_{\mathbb R}\otimes M_{\mathbb R},L_{\mathbb R})$ formé des formes $\varepsilon$-sym. \uline{non dégénérées}, et \inserted{sont \uline{donc} \uline{en nb fini}.}

\uline{Remarque 2} Les prépolarisations maximales de $\mathcal M^i$ ou $\mathcal M_{\mathbb R}^i$ sont en corr. biunivoque. De $=$ donc pour $\mathcal M$ et $\mathcal M_{\mathbb R}$. Si $P,P_{\mathbb R}$ se correspondent, alors $P$ est une polarisation sss $P_{\mathbb R}$ en est une. Si $\mathcal M$ est à engendrement fini, \deleted{\unsure{on trouve|c) Prouver}} l'ens des polarisations possibles de $\mathcal M$ est fini.

On est ainsi ramené à étudier les polarisations des $\otimes$-catégories \uline{sur} $\mathbb R$. De plus

\uline{Remarque 3} Soit $\mathcal M=\varinjlim\mathcal M_{\unsure{(\alpha)|\alpha}}$\deleted{\unsure{cat}} \inserted{sous-$\otimes$-catégories famille filtrante croissante}, alors l'ens $P(\mathcal M)$ des polarisations de $\mathcal M$ est donné par $P(\mathcal M)=\varprojlim_\alpha\mathcal M(\alpha)$ \deleted{\unsure{donc} $P(\mathcal M)$} Donc si les $\mathcal M(\alpha)$ \deleted{sont} de t.f., on

% archive sheet number: 7
